\documentclass[11pt]{article}
\usepackage[T1]{fontenc}
\usepackage{lmodern,amsmath,amssymb,amsthm,mathtools}
\usepackage[margin=30mm]{geometry}
\usepackage{microtype}
\usepackage{xcolor}
\usepackage[colorlinks=true,linkcolor=blue!45!black,citecolor=blue!45!black,urlcolor=blue!45!black]{hyperref}
\setlength{\emergencystretch}{2em}
\allowdisplaybreaks[2]
\numberwithin{equation}{section}
\newtheorem{theorem}{Theorem}[section]
\newtheorem{lemma}[theorem]{Lemma}
\newtheorem{proposition}[theorem]{Proposition}
\newtheorem{corollary}[theorem]{Corollary}
\theoremstyle{remark}
\newtheorem{remark}[theorem]{Remark}
\newcommand{\D}{\mathbb D}
\newcommand{\T}{\mathbb T}
\newcommand{\Q}{\mathbb Q}
\newcommand{\R}{\mathbb R}
\newcommand{\C}{\mathbb C}
\newcommand{\norm}[1]{\lVert #1\rVert}
\newcommand{\coef}[1]{[z^{#1}]}
\newcommand{\area}{\,dm}
\newcommand{\Rea}{\operatorname{Re}}
\title{Zeros and sharp coefficient bounds with a prescribed constant term in \(H^p\)}
\author{Research manuscript prepared for Keio}
\date{28 November 2026}
\hypersetup{pdftitle={Zeros and sharp coefficient bounds with a prescribed constant term in Hardy spaces},pdfauthor={Research manuscript prepared for Keio}}
\begin{document}
\maketitle

\begin{abstract}
We answer both questions posed by Brevig, Grepstad and Instanes
in Section~6 of \emph{F.~Wiener's Trick and an Extremal Problem for \(H^p\)}.
For \(0<p<1\), every extremal for their fixed-constant-term coefficient problem
has at most one disk zero, counted with multiplicity, and has no singular inner
factor. For every degree \(k\), the sharp value agrees with the first-coefficient
value whenever \(t\ge (1+p^2/(64k))^{-1/p}\), and the extremal is explicit and
unique for the real-part objective. At \(k=2\) we determine the entire
equality set, with exact lower endpoint
\((1+(p/(2-p))^2)^{-1/p}\).
For fixed \(p\), the width of the equality interval is of order
\(k^{-(2-2p)/(2-p)}\) as \(k\to\infty\), and we determine its leading
constant when \(2/p\) is an integer.
The proofs use perturbations along nested inner divisors and a coefficient
estimate that compares the linear gains from an inner factor with its cost
in the outer \(H^2\) norm. An arithmetic application gives an unconditional
three-value criterion equivalent to the irrationality of Euler's constant.
\end{abstract}

\section{The two fixed-constant-term questions}

Write \(dm\) for normalized arc measure on \(\T\). For \(0<p<1\) we use the
Hardy quasinorm \(\norm{f}_p=(\int_{\T}|f|^p\area)^{1/p}\), with radial boundary
values. Set
\begin{equation}\label{eq:phi}
 \Phi_k(p,t)=\sup\bigl\{\Rea\coef{k}f:\ f\in H^p,\ \norm{f}_p\le1,\ f(0)=t\bigr\},
 \qquad k\ge1,\quad 0\le t\le1 .
\end{equation}
A rotation \(z\mapsto e^{i\theta}z\) preserves feasibility and puts
a nonzero \(k\)-th coefficient on the positive real axis. Thus the sharp
bound for the modulus in this problem is the same as \eqref{eq:phi}.
A normal families argument gives attainment. Any maximizer has norm one.
Indeed, with norm slack a sufficiently small positive multiple of \(z^k\)
raises the objective and preserves feasibility, by
\(\norm{f+g}_p^p\le\norm{f}_p^p+\norm{g}_p^p\).

Brevig, Grepstad and Instanes asked in their 2020 preprint and in
\cite[Section~6]{BGI} whether the
extremals in \eqref{eq:phi} have at most one zero in \(\D\), and whether
\(\Phi_k(p,t)=\Phi_1(p,t)\) on an interval ending at \(1\) for each
\(k\ge2\) and \(0<p<1\). They gave numerical evidence for the first
question at \(k=2\) and numerical and analytic evidence for the second
at \(k=2\). Both questions also appear in Instanes's October 2024
lecture \cite[slides 25--26]{Instanes}. The fixed \(f(0)=t\)
condition matters. In particular, when \(t=0\) every feasible
function already vanishes at \(0\).

\begin{theorem}\label{thm:main}
Let \(0<p<1\) and \(k\ge1\).
\begin{enumerate}
\item For every \(0\le t\le1\), the inner factor of every extremal in
\eqref{eq:phi} is a constant or a Blaschke factor of degree one.
\item Put \(q=2/p\) and \(S=t^{-p}-1\) for \(t>0\).
If \(S\le p^2/(64k)\), then
\[
  \Phi_k(p,t)=q t\sqrt S=\Phi_1(p,t).
\]
The unique extremal for \eqref{eq:phi} is
\(f(z)=t(1+\sqrt S\,z^k)^q\).
\item For \(k=2\), one has \(\Phi_2(p,t)=\Phi_1(p,t)\) if and only if
\[
  t\ge t_*(p):=\bigl(1+(p/(2-p))^2\bigr)^{-1/p}.
\]
The same value and uniqueness formula from part~2 holds throughout
this sharp interval.
\end{enumerate}
\end{theorem}

The corpus's arbitrary-degree manuscript \cite{CorpusHardy} concerns
the unrestricted norm of the coefficient functional. Its polynomial
normalization and analysis of boundary zeros motivate the corresponding
questions on the affine slice \eqref{eq:phi}. Our proofs use the original
slice in \cite{BGI}. The two main steps are a dimension bound for
real boundary perturbations that preserve affine constraints, and a sharp
estimate for a coefficient under a small outer \(H^2\) budget.
The arithmetic section takes the cancellation approach of the corpus's
Sondow manuscript \cite{CorpusSondow} to a different family of quantities.

\section{A zero bound under affine constraints}

We first formulate the auxiliary question for arbitrary bounded linear
constraints. Linear maps and functionals in this section are complex
linear. Constants are allowed among the inner functions.

\begin{lemma}\label{lem:prefix}
Let \(0<p<1\), \(f\ne0\) in \(H^p\), and suppose the inner factor of
\(f\) is divisible by \(J_1\cdots J_d\) with each \(J_j\) nonconstant
inner. There is a real vector space of dimension \(2d+1\) of bounded
real boundary functions \(h\) for which \(fh\) is the boundary
function of an element of \(H^p\).
\end{lemma}
\begin{proof}
Put \(U_j=J_1\cdots J_j\), and take
\[
 h=c_0+\sum_{j=1}^d(c_j U_j+\overline{c_j}\,\overline{U_j})
 \quad\text{on }\T,\qquad c_0\in\R,\quad c_j\in\C .
\]
The analytic representative for \(fh\) is
\(c_0 f+\sum_j(c_j fU_j+\overline{c_j}f/U_j)\), which belongs to \(H^p\)
by inner divisibility. To check dimension, \(1,U_1,\ldots,U_d\)
are linearly independent over \(\C\). In a linear relation, a nonzero
constant term would give a bounded analytic inverse to \(J_1\).
Such an inverse would have boundary modulus one and hence modulus
at most one in the disk, which is impossible for the inverse of a
nonconstant inner function. After the constant term vanishes, cancel
\(J_1\) and iterate. Finally, a bounded analytic function with
zero real part almost everywhere on \(\T\) is an imaginary constant.
Applying this to \(c_0/2+\sum_j c_j U_j\) proves real linear independence
of the displayed parameters.
\end{proof}

\begin{theorem}\label{thm:affine}
Let \(A:H^p\to\C^m\) and \(L:H^p\to\C\) be bounded linear maps.
Suppose \(L\) is nonzero on \(\ker A\). Every maximizer of \(\Rea L\)
over
\(\{f:\norm{f}_p\le1,\ Af=w\}\)
has a finite Blaschke inner factor of degree at most \(m\).
For \(m\ge1\), every nonzero norm minimizer of \(Af=w\) has a finite
Blaschke inner factor of degree at most \(m-1\).
Both bounds are attained for suitable choices of the constraints.
\end{theorem}
\begin{proof}
Choose \(v\in\ker A\) with \(\Rea L(v)>0\).
For the maximization problem, norm slack permits addition of
\(\delta v\) for small positive \(\delta\).
Consequently an extremal \(f\) has norm one.

If the inner factor of \(f\) has at least \(m+1\) disk zeros counted
with multiplicity, take \(d=m+1\) elementary factors.
If it has a nonconstant singular factor \(S_\mu\), take \(d=m+1\) copies
of \(S_{\mu/d}\). In either case Lemma~\ref{lem:prefix} gives \(2m+3\)
real parameters for \(h\). Impose
\[
 A(fh)=0,\qquad \int_{\T} h|f|^p\area=0 ,
\]
at most \(2m+1\) real conditions. A nonzero solution \(h\) remains.
Since \(|f|>0\) almost everywhere and the weighted mean is zero, \(h\)
is not almost everywhere constant for the probability measure
\(|f|^p dm\). For small \(\epsilon>0\), \(1\pm\epsilon h>0\) on
\(\T\) up to a null set. Strict concavity of \(x^p\) gives
\[
  \norm{f\pm\epsilon fh}_p^p
   =\int_{\T}(1\pm\epsilon h)^p |f|^p\area<1.
\]
Both variations satisfy \(Af=w\), and at least one has objective at
least \(\Rea L(f)\). The resulting norm slack contradicts maximality.

For the minimization problem, factor out \(d=m\) nonconstant inner
functions if possible, and impose only \(A(fh)=0\). This again leaves
a nonzero \(h\). The average of the \(p\)-th powers of the two perturbed norms
is strictly below \(\norm{f}_p^p\) by pointwise strict concavity on the
set where \(h\ne0\). Thus at least one feasible perturbation improves
the norm, a contradiction.

Sharpness in the maximization statement follows by imposing the
vanishing of the first \(m\) coefficients and maximizing
\(\Rea\coef{m}f\). Division by \(z^m\) preserves the norm and reduces
the problem to point evaluation at the origin. Its unique extremal
is \(1\), by subharmonicity and the strict mean-value inequality for
a nonconstant analytic function. In the minimization statement
prescribe the first \(m-1\) coefficients to be zero and
\(\coef{m-1}f=1\). Division by \(z^{m-1}\) gives the same sharpness
argument.
\end{proof}

Take \(A(f)=f(0)\) and \(L(f)=\coef{k}f\) to obtain
Theorem~\ref{thm:main}(1). The real perturbation is related to the
de Leeuw--Rudin construction in \(H^1\) \cite{dLR}, and to
Brevig--S{\o}vik's unconstrained result for \(0<p<1\) \cite{Sovik}.
Dyakonov's dimension and rank argument for the homogeneous punctured
\(H^1\) ball is a close precedent \cite[Theorem~1.1 and Section~3]{Dyakonov}.
Here the nested divisors give the perturbations in \(H^p\), including for
singular factors. Strict concavity makes the norm strictly decrease
while the affine constraints stay fixed.

\section{Boundary zeros}

For \eqref{eq:phi}, the classical extremal form is
\begin{equation}\label{eq:kab}
  f=c\,B P^q,\qquad
  P(z)=\prod_{j=1}^k(1-\overline{\lambda_j}z),
  \quad |\lambda_j|\le1 ,
\end{equation}
where \(B\) is a finite Blaschke product whose zeros, counted with
multiplicity, are among the \(\lambda_j\) in \(\D\). Factors with
\(\lambda_j=0\) in \(P\) are interpreted as \(1\).
This is the form attributed to Kabaila in \cite[Lemma~4]{BGI}. The
power is taken with \(P^q(0)=1\). We use this classical input for the
boundary statement and for the converse in Section~\ref{sec:k2}.

\begin{proposition}\label{prop:boundary}
For an extremal in \eqref{eq:phi}, let \(N\) be the number of disk zeros
of \(f\) and \(M\) the number of circle zeros of \(P\) in \eqref{eq:kab},
both counted with multiplicity. Then \(2N+M\le2\).
\end{proposition}
\begin{proof}
For a disk zero \(a\) of multiplicity \(r\), include in the space of
boundary multipliers all functions
\[
  \frac{c_j}{(z-a)^j}
       +\overline{c_j}\frac{z^j}{(1-\overline a z)^j},
      \qquad c_j\in\C,\quad 1\le j\le r .
\]
For a zero \(\zeta\in\T\) of \(P\) of multiplicity \(r\), include the
real span of
\[
  \left(i\frac{\zeta+z}{\zeta-z}\right)^j,\qquad 1\le j\le r .
\]
All these functions are real on \(\T\) away from poles.
Their principal parts at the chosen poles give real dimension \(2N+M\),
and adjoining constants adds one dimension. The poles of an interior
multiplier cancel in \(fh\). At a boundary pole of order at most \(r\),
the factor \(P^q\) vanishes to order \(qr\) in modulus and \(q>2\).
Thus \(fh\) is analytic and bounded in \(\D\).
Also \(h\in L^2(|f|^p dm)\), since \(|f|^p\) vanishes to order \(2r\)
at that boundary point.

If \(2N+M>2\), impose \((fh)(0)=0\) and
\(\int_\T h|f|^p\area=0\), leaving a nonzero \(h\).
For real \(x\) the function
\(\bigl(|1+x|^p-1-px\bigr)/x^2\) is bounded, with limiting value
\(p(p-1)/2\) at zero. Dominated convergence therefore yields
\[
  \norm{f+\epsilon fh}_p^p
   =1+\frac{p(p-1)}2\epsilon^2\int_\T h^2|f|^p\area+o(\epsilon^2).
\]
The second-order coefficient is negative. The two signs give
norm slack, and one of them retains at least the old objective.
Adding a small positive multiple of \(z^k\) is a contradiction.
\end{proof}

Thus a fixed-constant-term extremal with one disk zero has an outer
polynomial without circle zeros. When there is no disk zero, the argument
allows at most two circle zeros.

\section{Exact symmetry near \texorpdfstring{\(t=1\)}{t=1}}

We work with every feasible function and use no extremal-form theorem in
this section. Let \(0<t\le1\) and \(q=2/p\). Factor \(f=IO\) as inner times
outer and normalize the phases so \(I(0)=\beta>0\).
Then \(O(0)=t/\beta\) and
\[
 P=(O/O(0))^{1/q},\quad P(0)=1,\quad
 f/t=(I/\beta)P^q,\quad
 1+E:=\norm{P}_2^2\le(\beta/t)^p .
\]
In particular \(t\le\beta\le1\), and \(S=t^{-p}-1\) satisfies
\begin{equation}\label{eq:budget}
  E\le S,\qquad S-E\ge(1+S)(1-\beta^p).
\end{equation}

\begin{lemma}\label{lem:inner}
For an inner function with \(I(0)=\beta\in[0,1]\),
\(\lvert\coef{j}I\rvert\le1-\beta^2\) for each \(j\ge1\).
\end{lemma}
\begin{proof}
Average \(I\) over the \(j\)-th roots of unity. This bounded
analytic function has the form \(G(z^j)\) with
\(\norm{G}_\infty\le1\), \(G(0)=\beta\) and
\(G'(0)=\coef{j}I\). Apply Schwarz--Pick at the origin.
\end{proof}

\begin{proof}[Proof of Theorem~\ref{thm:main}(2)]
The case \(S=0\) follows from \eqref{eq:budget}. Otherwise write
\(b=\coef{k}P\) and \(R=\sum_{j=1}^{k-1}(\coef{j}P)z^j\), with the
empty sum allowed. Put \(v^2=\norm{R}_2^2\) and
\(x=\sqrt{kS}\). For \(j\ge2\), a coefficient of \(R^j\) has absolute
value at most \(v^2\norm{R}_\infty^{j-2}\) by integration on
\(\T\). Also \(\norm{R}_\infty\le\sqrt{k}v\le x\).
The order-\(k\) coefficient of \(P^q\) is \(qb\) plus terms coming
from \(R^j\) with \(2\le j\le k\). Since
\(\lvert\binom{q}{j}\rvert\le q^j\) and \(qx\le 1/4\), it follows that
\begin{equation}\label{eq:curv}
   \lvert\coef{k}P^q-qb\rvert\le\frac43 q^2 v^2.
\end{equation}
The same binomial expansion with the absolute coefficient sum of
\(R\) gives
\begin{equation}\label{eq:lone}
   \sum_{j=0}^{k-1}|\coef{j}P^q|
     \le \sum_{\ell\ge0}(qx)^\ell \le \frac43 .
\end{equation}
Here the expansions are Taylor identities to finite order. They
hold for real \(q\), and the bounds use the absolute values of
the binomial coefficients.

Lemma~\ref{lem:inner} and \eqref{eq:lone} bound the contribution of
\(I/\beta-1\) to \(\coef{k}f/t\) by
\(\tfrac43(1-\beta^2)/\beta\). Since
\(\beta\ge t=(1+S)^{-q/2}>1/2\) and
\(1-\beta^2=1-(\beta^p)^q\le q(1-\beta^p)\), this contribution
is at most \(4q(S-E)\) by \eqref{eq:budget}.
Moreover
\[
 \sqrt S-|b|\ge\frac{S-|b|^2}{2\sqrt S},\qquad
 v^2\le E-|b|^2.
\]
Combine these inequalities to obtain
\begin{align*}
 q\sqrt S-\frac{|\coef{k}f|}{t}
 &\ge \left(\frac{q}{2\sqrt S}-4q\right)(S-E)\\
 &\quad+\left(\frac{q}{2\sqrt S}-\frac43q^2\right)(E-|b|^2).
\end{align*}
Both coefficients are positive, since
\(S\le1/(16q^2k)\) and \(q>2\).
Equality forces \(E=|b|^2=S\), \(\beta=1\), and
\(P=1+bz^k\). The real-part objective requires \(b=\sqrt S\).
Conversely this \(P\) has a zero-free disk power \(tP^q\), with
\(\norm{tP^q}_p^p=t^p(1+S)=1\) and coefficient \(qt\sqrt S\).
The condition for degree \(k\) implies that for degree one,
which proves the equality of the two sharp values.
\end{proof}

\section{The exact quadratic threshold}\label{sec:k2}

We retain the notation \(\beta,E,S,q\) above. The following estimate
will compare the loss of \(H^2\) budget with the contribution of a
degree-one inner function. Suppose \(0<S\le1/(q-1)\). For \(0\le u\le S\) set
\[
  \beta_0(u)=\left(\frac{1+u}{1+S}\right)^{q/2}.
\]
For \(0<u\le S\), the function \((1+u)^{q/2}/\sqrt u\)
is decreasing on this interval, by logarithmic differentiation.
Consequently
\begin{equation}\label{eq:beta-min}
  \beta\ge\beta_0(E),\qquad
  \beta_0(u)\ge\sqrt{u/S}\ge\sqrt u.
\end{equation}
Because \(y+y^{-1}\) decreases on \(0<y\le1\),
\[
  -\frac{d}{du}\bigl(\beta_0(u)^{-1}-\beta_0(u)\bigr)
   =\frac{q}{2}\frac{\beta_0(u)+\beta_0(u)^{-1}}{1+u}
   \le\frac{q}{2\sqrt u}.
\]
Integration, with a limit at \(u=0\) if needed, gives
\begin{equation}\label{eq:beta-cost}
   \beta^{-1}-\beta\le q(\sqrt S-\sqrt E).
\end{equation}
The inequality is strict for \(E<S\), since \(S<1\).
In particular Lemma~\ref{lem:inner} gives
\(|\coef{1}f|/t\le q\sqrt E+\beta^{-1}-\beta\le q\sqrt S\).
The feasible function \(t(1+\sqrt S\,z)^q\) attains this bound.

Now set \(S_*=(q-1)^{-2}\). Take \(0<S\le S_*\) and an extremal
for degree two. By Theorem~\ref{thm:affine} its inner factor is either
constant or has degree one. Rotate to write the latter as
\(I(z)=(\beta+z)/(1+\beta z)\).
This rotation preserves the modulus of the coefficient.
Write \(a=\coef{1}P\), \(b=\coef{2}P\), \(r=|a|\), and
\(C=1-\beta^2\). The coefficient of \(f/t\) is
\begin{equation}\label{eq:coeff2}
  qb+\frac{q(q-1)}2a^2+\frac{q C}{\beta}a-C .
\end{equation}
The same expression with \(C=0\) applies to a constant inner factor.

For nonnegative \(A,J,C\) with \(AC>0\), completing the square in
\(\cos\theta\) gives
\[
  |Ae^{2i\theta}+Je^{i\theta}-C|^2
  \le (A+C)^2\left(1+\frac{J^2}{4AC}\right).
\]
Indeed the left side is
\((A+C)^2+J^2+2J(A-C)\cos\theta-4AC\cos^2\theta\).
In \eqref{eq:coeff2} use \(A=q(q-1)r^2/2\) and \(J=qCr/\beta\).
The factor in parentheses is at most \(1/\beta^2\), since
\(J^2/(4AC)=qC/(2(q-1)\beta^2)\).
The consequence \(\lvert Ae^{2i\theta}+Je^{i\theta}-C\rvert
\le (A+C)/\beta\) also holds when \(r=0\) or \(C=0\) by
direct inspection. Thus
\[
 \frac{|\coef{2}f|}{t}
   \le q|b|+\frac{q(q-1)r^2}{2\beta}+\beta^{-1}-\beta .
\]
From \eqref{eq:beta-min}, \(\beta\ge(q-1)\sqrt E\).
When \(E>0\), \( |b|\le\sqrt{E-r^2}\) implies
\[
 q|b|+\frac{q(q-1)r^2}{2\beta}
 \le q\left(\sqrt{E-r^2}+\frac{r^2}{2\sqrt E}\right)
 \le q\sqrt E.
\]
When \(E=0\) the same bound holds because \(a=b=0\).
By \eqref{eq:beta-cost}, the desired bound \(q\sqrt S\) follows.
Equality implies \(E=S\), \(\beta=1\), and \(r=0\) by the strict
concavity of the square root. Equality in the estimate for \(b\)
then forces \(P=1+bz^2\) with \(|b|=\sqrt S\). The unique
positive-real maximizer follows, including \(S=0\).

It remains to prove strict inequality below \(t_*\).
If equality holds, composing any first-coefficient extremal \(g\)
with \(z^2\) gives a degree-two extremal with the same value.
Theorem~\ref{thm:affine} forces \(g\) to be outer.
Indeed a single disk zero of \(g\) gives two disk zeros counted
with multiplicity after composition. A singular inner factor has
already been excluded for \(g\) itself. At \(t=0\) an outer \(g\) is
impossible. For \(t>0\), the classical form \eqref{eq:kab} at \(k=1\)
then gives \(g(z)=t(1+\sqrt S\,z)^q\), with \(\sqrt S\le1\).
It has first coefficient \(qt\sqrt S\).
If \(S_*<S\le1\), take a small positive \(a\) and
\[
  P_a(z)=1+az+\sqrt{S-a^2}\,z^2.
\]
For \(S<1\), the polynomial is zero-free on the closed disk for
sufficiently small \(a\) by uniform perturbation of \(1+\sqrt S\,z^2\).
For \(S=1\), it has conjugate roots for sufficiently small \(a\),
each of modulus \((\sqrt{1-a^2})^{-1/2}>1\).
The feasible \(tP_a^q\) has second coefficient divided by \(t\)
equal to
\[
   q\sqrt{S-a^2}+\frac{q(q-1)}2 a^2>q\sqrt S
\]
for small \(a>0\), since the derivative with respect to \(a^2\) at
zero is positive. This contradicts equality and proves
Theorem~\ref{thm:main}(3).

\begin{remark}
Section~6 of \cite{BGI} prints the \(k=2\) proposed endpoint with a
positive exponent \(1/p\). That value is outside \([0,1]\).
The endpoint proved here has exponent \(-1/p\).
For each \(k\ge3\) the quadratic endpoint does not suffice.
At \(S=S_*\), use \(P=1+a(z+\cdots+z^{k-1})+
\sqrt{S-(k-1)a^2}\,z^k\) for small \(a>0\).
This is zero-free on the closed disk by continuity. In the
\(k\)-th coefficient of \(P^q\) the order-\(a^2\) terms cancel
and the order-\(a^3\) term is
\(\binom{q}{3}\binom{k-1}{2}a^3>0\).
Thus \(\Phi_k(p,t_*)>\Phi_1(p,t_*)\) for every \(k\ge3\).
\end{remark}

\section{The length of the symmetry interval}

For fixed \(p\) and \(k\ge2\), define \(w_k(p)\) as the supremum of
\(W\ge0\) such that \(\Phi_k(p,t)=\Phi_1(p,t)\) for all \(t\) with
\(0\le t^{-p}-1\le W\).

\begin{theorem}\label{thm:scale}
For each \(0<p<1\), there are constants \(c_p>0\) and \(C_p'>0\) such that
\[
 c_p k^{-(2-2p)/(2-p)} \le w_k(p)\le
 C_p' k^{-(2-2p)/(2-p)}
\]
for all sufficiently large \(k\). On the interval
\(0\le t^{-p}-1\le c_p k^{-(2-2p)/(2-p)}\),
the unique extremal for the real-part objective is
\(t(1+\sqrt{t^{-p}-1}\,z^k)^{2/p}\).
\end{theorem}

To prove the lower bound with a real noninteger power, we need to
control the argument of a normalized analytic function on the set
where it is close to one.

\begin{lemma}\label{lem:L1}
Let \(q>2\). Suppose \(1+x\) is analytic and zero-free on a neighborhood
of the closed disk, and \(x(0)=0\).
Use the analytic power with \((1+x)^q(0)=1\).
There is a constant \(K_q\) independent of \(x\) such that
\[
  \int_\T |(1+x)^q-1-qx|\area
     \le K_q\left(\int_\T |x|^2\area+\int_\T |x|^q\area\right).
\]
\end{lemma}
\begin{proof}
Put \(G=\log(1+x)-x\), with \(\log(1+x)(0)=0\).
Since \(\log(1+y)\le y\) for \(y>-1\),
\[
 \Rea G=\tfrac12\log(1+2\Rea x+|x|^2)-\Rea x\le |x|^2/2 .
\]
Let \(J\) be the analytic function whose real part is the Poisson
extension of \(|x|^2/2\) on \(\T\) and whose imaginary part vanishes at
zero. Both \(J\) and \(J-G\) have nonnegative real part and value
\(v=\int_\T |x|^2\area/2\) at zero. They have continuous boundary
values, since \(x\) is analytic near \(\T\).
For any such function \(U\) and \(T>0\), the real part of
\(U/(T+U)\) is nonnegative and is at least \(1/2\) wherever
\(|U|\ge T\). The mean-value property gives
\[
 m\{|U|\ge T\}\le \frac{2v}{T+v}.
\]
Applying this to \(J,J-G\) with \(T=1/2\) shows
\(m\{|\operatorname{Im}G|>1\}\le8v\).

On \(\{|x|\le1/2,\ |\operatorname{Im}G|\le1\}\), the imaginary part of
the normalized logarithm is bounded by \(3/2\), so this logarithm is
the principal one there. Taylor's formula gives a pointwise bound
\(K_q|x|^2\) for the remainder. On the rest of \(\{|x|\le1/2\}\) the
remainder is bounded by a constant depending only on \(q\), since
\(|(1+x)^q|=|1+x|^q\). That set has measure at most \(8v\).
On \(\{|x|>1/2\}\) the remainder is at most \(K_q(|x|^2+|x|^q)\),
again by the same identity for the modulus. Integration proves the
lemma.
\end{proof}

\begin{proof}[Proof of Theorem~\ref{thm:scale}]
Use \(q=2/p\) and \(S=t^{-p}-1\). For the lower bound, let \(f\)
be any extremal, with the factorization and notation from Section~4.
Put \(b=\coef{k}P\), \(\Delta=E-|b|^2\), and \(\varepsilon=S-E\).
We may restrict to sufficiently small \(S\) depending on \(q\),
so that \(\beta\ge1/2\) and \(|b|\le1/2\).
Let \(P_0=1+bz^k\) and \(\rho=1-1/k\).
Apply Lemma~\ref{lem:L1} to
\[
  1+x=\frac{P(\rho z)}{P_0(\rho z)} .
\]
Its numerator and denominator are nonzero on a neighborhood of the
closed disk. Parseval and Cauchy--Schwarz give
\(\norm{x}_2\le2\sqrt{\Delta}\) and
\(\norm{x}_\infty\le2\sqrt{k\Delta}\) on \(\T\).
Multiplying the remainder in Lemma~\ref{lem:L1} by \(P_0(\rho z)^q\)
therefore gives an \(L^1\) norm at most
\[
 O_q\bigl(\Delta(1+(k\Delta)^{(q-2)/2})\bigr).
\]
The Taylor coefficients up to order \(k\) of
\(P_0^q+qP_0^{q-1}(P-P_0)\) are those of \(1+q(P-1)\).
Thus, by the coefficient integral and \(\rho^{-k}\le4\),
\begin{equation}\label{eq:remainder}
 |\coef{j}(P^q-1-q(P-1))|
   \ll_q \Delta(1+(k\Delta)^{(q-2)/2}),\qquad 0\le j\le k .
\end{equation}

Theorem~\ref{thm:affine} gives degree at most one for \(I\).
If \(I\) has degree one, the absolute sum of its nonconstant
Taylor coefficients divided by \(\beta\) is
\((1+\beta)/\beta\le3\), by the geometric series for a Blaschke
factor. It is zero if \(I\) is constant.
The \(H^2\) norm of \(I/\beta-1\) is
\(\sqrt{1-\beta^2}/\beta\le2\sqrt{q\varepsilon}\) by
\eqref{eq:budget}. Its \(k\)-th coefficient has modulus at most
\(2q\varepsilon\) by Lemma~\ref{lem:inner}.
Expand \((I/\beta)P^q\) using \eqref{eq:remainder}.
The \(k\)-th coefficient of \((I/\beta-1)(P-1)\) has modulus at most
\(2\sqrt{q\varepsilon\Delta}\) by Cauchy--Schwarz.
The sum of the remaining nonlinear terms is bounded using the
absolute coefficient sum of \(I/\beta\), which is at most four.
Consequently
\begin{equation}\label{eq:scaled-error}
 \frac{|\coef{k}f|}{t}
   \le q|b|+O_q\bigl(\varepsilon+\sqrt{\varepsilon\Delta}
                 +\Delta(1+(k\Delta)^{(q-2)/2})\bigr).
\end{equation}
The loss \(q(\sqrt S-|b|)\) is at least
\(q(\varepsilon+\Delta)/(2\sqrt S)\).
For \(S\le c_p k^{-(q-2)/(q-1)}\), one has
\[
 \sqrt S\,(kS)^{(q-2)/2}\le c_p^{(q-1)/2}.
\]
Together with \(\sqrt{\varepsilon\Delta}\le(\varepsilon+\Delta)/2\),
this shows that the loss strictly dominates the error in
\eqref{eq:scaled-error} whenever \(\varepsilon+\Delta>0\) if \(c_p>0\)
is small enough. The equality case \(\varepsilon=\Delta=0\)
gives the stated monomial factor and uniqueness, as in Section~4.
For large \(k\), this range lies within the first-coefficient
interval of Theorem~\ref{thm:main}(2).

For the upper bound, let \(r=1-1/k\) and take the norm-one function
\[
  G(z)=(1-r^2)^{1/p}(1-rz)^{-2/p}.
\]
Its coefficient satisfies
\[
  \coef{k-1}G
      =(1-r^2)^{1/p}r^{k-1}
       \frac{\Gamma(k-1+2/p)}{\Gamma(2/p)\Gamma(k)}
      \gg_p k^{1/p-1}.
\]
The function \(f=t+(1-t^p)^{1/p}zG\) is feasible by the
\(p\)-triangle inequality. Thus
\(\Phi_k(p,t)\gg_p t S^{1/p} k^{1/p-1}\).
At \(S=C_p' k^{-(q-2)/(q-1)}\) with \(C_p'\) sufficiently large,
this strictly exceeds the first-coefficient value \(qt\sqrt S\)
for all large \(k\). This proves the claimed upper bound.
\end{proof}

\begin{theorem}\label{thm:phaseconstant}
If \(q=2/p\) is an integer at least three, put
\[
  \tau=q/2-1,\qquad
  A_p=\lim_{j\to\infty}\frac{C(j,p)}{j^\tau},\qquad
  M_q=\max_{0\le x\le1}(1+x)^{q/2}(1-x)^{q/2-1}.
\]
Here \(C(j,p)=\sup_{\norm{f}_p\le1}|\coef{j}f|\).
The limit \(A_p\) exists and is positive, and the maximum defining
\(M_q\) is attained at \(x=1/(q-1)\). Moreover,
\[
  w_k(p)\sim
   \left(\frac{q}{A_p M_q}\right)^{2/(q-1)}
       k^{-(q-2)/(q-1)} .
\]
\end{theorem}
\begin{proof}
For clarity the short argument for the limit, also used in
\cite{CorpusLimits}, is included. Projection onto the coefficients
with indices divisible by \(m\), followed by replacing \(z^m\) with
\(z\), gives \(C(mj,p)\le m^\tau C(j,p)\) by the \(p\)-triangle
inequality. Also \(C(j,p)\) is nondecreasing by multiplication by
\(z\). For each fixed \(j\), compare \(C(n,p)\) to
\(C(j\lceil n/j\rceil,p)\) to obtain
\(\limsup_{n\to\infty} C(n,p)/n^\tau\le C(j,p)/j^\tau\).
Hence the limit is the infimum of these ratios and is positive
by the kernel test used above. The maximum in \(M_q\) follows by
logarithmic differentiation on \((0,1)\) and checking the endpoints.

Let \(\alpha=(q-2)/(q-1)\) and fix \(D>0\). We first work uniformly
over \(0<S\le D k^{-\alpha}\).
Write \(R=P-(1+b z^k)\), so \(\norm{R}_2^2=\Delta\).
Since \(q\) is an integer, through order \(k\) the powers of
\(P-1\) of order \(m\ge2\) agree with \(R^m\).
For \(2\le m<q\), all coefficients through order \(k\) of \(R^m\)
are bounded in modulus by \(\Delta(k\Delta)^{(m-2)/2}\).
This follows by truncating \(R\) to order \(k\) and using its
\(L^2\) and \(L^\infty\) norms as in \eqref{eq:curv}.
The estimates for the inner factor from the proof of
Theorem~\ref{thm:scale} therefore give
\[
 \frac{|\coef{k}f|}{t}
 \le q|b|+\beta^{-1} C(k,p)\Delta^{q/2}
       +o(1)(\sqrt S-|b|).
\]
Indeed the top term is \(\coef{k}((I/\beta)R^q)\), bounded as
displayed since \(\norm{R^q}_p=\norm{R}_2^q\). Every lower term,
after division by \(\sqrt S-|b|\), is bounded by a constant times
\(\sqrt S(1+\sum_{m=2}^{q-1}(kS)^{(m-2)/2})=o(1)\),
uniformly in this range. When \(\sqrt S=|b|\) the lower terms
vanish. Also \(\beta^{-1}=1+o(1)\) uniformly.

The definition of \(M_q\) gives
\[
  \Delta^{q/2}
   \le(S-|b|^2)^{q/2}
   \le M_q S^{(q-1)/2}(\sqrt S-|b|).
\]
Thus for every fixed \(D<(q/(A_p M_q))^{2/(q-1)}\) and sufficiently
large \(k\), the monomial factor is the unique maximizer
throughout \(0\le S\le D k^{-\alpha}\).
For the reverse inequality, choose
\(D>(q/(A_p M_q))^{2/(q-1)}\) and set \(S=D k^{-\alpha}\).
Put \(b=\sqrt S/(q-1)\), \(\Delta_0=S-b^2\), and choose a norm-one
extremal \(G\) with \(\coef{k-1}G=C(k-1,p)\).
Then
\[
    f=t\bigl((1+bz^k)^q+\Delta_0^{q/2} zG\bigr)
\]
is feasible by the \(p\)-triangle inequality.
Its \(k\)-th coefficient divided by \(t\) is
\(qb+C(k-1,p)\Delta_0^{q/2}\). This exceeds \(q\sqrt S\)
for all large \(k\), since
\(C(k-1,p) M_q S^{(q-1)/2}\to A_p M_q D^{(q-1)/2}>q\).
The two bounds on \(D\) give the asymptotic.
\end{proof}

\section{A three-value criterion for Euler's constant}

For the unrestricted coefficient norm \(C(2,p)\) just defined,
Brevig--Saksman \cite[Theorem~1]{BS} give, for \(s=2/p>2\),
\[
   F(s):=C(2,2/s)=(s-1)^{s-1}s^{2-s}.
\]
The following gives a two-way arithmetic test for the original
Euler--Mascheroni constant
\(\gamma=\lim_{n\to\infty}(\sum_{j=1}^n j^{-1}-\log n)\).

\begin{proposition}\label{prop:gamma}
A real number \(s>3\) is rational if and only if all three numbers
\(F(s),F(s+1),F((s+1)/2)\) are algebraic.
In particular, \(\gamma\) is irrational if and only if at least one of
\[
  C(2,\gamma),\quad
  C\left(2,\frac{2\gamma}{2+\gamma}\right),\quad
  C\left(2,\frac{4\gamma}{2+\gamma}\right)
\]
is transcendental.
\end{proposition}
\begin{proof}
Direct multiplication cancels the variable powers and gives
\[
  \frac{F(s)F(s+1)}{F((s+1)/2)^2}
      =\left(\frac{2s}{s+1}\right)^2.
\]
If all three values are algebraic, the positive square root shows
that \(r=2s/(s+1)\) is algebraic, and
\(s=r/(2-r)\) is algebraic. An algebraic irrational \(s\) is
impossible in this case. Indeed
\(F(s)=s(1-1/s)^{s-1}\) is then transcendental by
Gelfond--Schneider \cite{Baker}.
Conversely if \(s\) is rational, the three expressions are products
of positive rational bases to rational powers and hence algebraic.
For \(\gamma\) use \(s=2/\gamma>3\), as \(1/2<\gamma<2/3\), and substitute.
\end{proof}

The cancellation identity isolates a further arithmetic step.
An independent transcendence proof for one of the three values
would decide irrationality. The identity and the analytic extremal
results above do not supply that proof. Sondow's criterion \cite{Sondow}
and the recent recurrence approach of \cite{Kalisch} likewise
leave the irrationality question open.

\section{Relation to prior work}

The source problem for Theorem~\ref{thm:main}(1)--(2) is
\cite[Section~6, Questions~1--2]{BGI}. Both answers are for the original
\(\Phi_k(p,t)\) on the Hardy unit ball. The first covers every
\(0\le t\le1\). The second covers every requested exponent and degree
and supplies an explicit interval, as the question asks, with a sharp
endpoint for \(k=2\).
The affine real-dimension argument uses a mechanism present in
\cite{dLR,Dyakonov,Sovik}. The explicit strict deficit in Section~4
also handles all feasible functions and all real \(0<p<1\), without
a classification of competing stationary points.
The closest sources located were the original article, the 2024
lecture, the result on the punctured \(H^1\) ball, and the unconstrained
outer extremal theorem. These sources treat the respective antecedents
described above. No earlier answer to the two fixed-constant-term
questions was identified in the sources checked. This is a limited
literature finding and priority remains subject to independent review.

\begin{thebibliography}{12}

\bibitem{Baker}
A.~Baker, \emph{Transcendental Number Theory}, Cambridge University
Press, 1975.

\bibitem{BGI}
O.~F.~Brevig, S.~Grepstad and S.~M.~Instanes,
F.~Wiener's trick and an extremal problem for \(H^p\),
\emph{Comput. Methods Funct. Theory} \textbf{23} (2023), 697--722.
\href{https://doi.org/10.1007/s40315-022-00469-x}{DOI}.
\href{https://arxiv.org/abs/2011.07790v1}{2020 preprint}.

\bibitem{BS}
O.~F.~Brevig and E.~Saksman,
Coefficient estimates for \(H^p\) spaces with \(0<p<1\),
\emph{Proc. Amer. Math. Soc.} \textbf{148} (2020), 3911--3924.
\href{https://arxiv.org/abs/1905.09547}{Author manuscript}.

\bibitem{Sovik}
O.~F.~Brevig and {\O}.~S{\o}vik,
Extremals for linear functionals on \(H^p\) are outer when \(0<p<1\),
preprint (2026).
\href{https://arxiv.org/abs/2610.21788}{arXiv:2610.21788}.

\bibitem{dLR}
K.~de Leeuw and W.~Rudin,
Extreme points and extremum problems in \(H_1\),
\emph{Pacific J. Math.} \textbf{8} (1958), 467--485.
\href{https://doi.org/10.2140/pjm.1958.8.467}{DOI}.

\bibitem{Duren}
P.~L.~Duren, \emph{Theory of \(H^p\) Spaces},
Pure and Applied Mathematics, vol.~38, Academic Press, 1970.

\bibitem{Dyakonov}
K.~M.~Dyakonov, Nearly outer functions as extreme points in punctured
Hardy spaces, \emph{Adv. Math.} \textbf{401} (2022), 108330.
\href{https://doi.org/10.1016/j.aim.2022.108330}{DOI}.

\bibitem{Instanes}
S.~M.~Instanes, \emph{An extremal problem for \(H^p\)},
lecture slides, 2 October 2024.
\href{https://verso.mat.uam.es/~dragan.vukotic/respub/sem_2024-10-02_Instanes-slides.pdf}{Slides}.

\bibitem{Kalisch}
R.~Kalisch, T.~Monderer, E.~Leibtag, S.~Weinbaum, M.~Shalyt and
I.~Kaminer, An infinite family of Ap\'ery-like recurrences for Euler's
constant, preprint (2026).
\href{https://arxiv.org/abs/2610.13919}{arXiv:2610.13919}.

\bibitem{CorpusHardy}
\emph{Verifiable Mathematical Reasoning Corpus}, item 279,
Sharp Hardy coefficients in arbitrary degree,
snapshot \href{https://github.com/QEEDDD/Verifiable-Mathematical-Reasoning-Corpus/tree/9f8547bfd3b8f6560f20c1b63028dc87a07a2e33}{\texttt{9f8547b}} (2026).

\bibitem{CorpusLimits}
\emph{Verifiable Mathematical Reasoning Corpus}, item 280,
Hardy coefficient limits and Dirichlet maximal order,
snapshot \href{https://github.com/QEEDDD/Verifiable-Mathematical-Reasoning-Corpus/tree/9f8547bfd3b8f6560f20c1b63028dc87a07a2e33}{\texttt{9f8547b}} (2026).

\bibitem{CorpusSondow}
\emph{Verifiable Mathematical Reasoning Corpus}, item 100,
Sondow integrals and three-grid Cantor series,
snapshot \href{https://github.com/QEEDDD/Verifiable-Mathematical-Reasoning-Corpus/tree/9f8547bfd3b8f6560f20c1b63028dc87a07a2e33}{\texttt{9f8547b}} (2026).

\bibitem{Sondow}
J.~Sondow, Criteria for irrationality of Euler's constant,
\emph{Proc. Amer. Math. Soc.} \textbf{131} (2003), 3335--3344.
\href{https://arxiv.org/abs/math/0209070}{Author manuscript}.

\end{thebibliography}
\end{document}
